\pagestyle{plain}

\chapter{CAPÍTULO 6: Design Participativo de Sistema de Governança Bioeconômica para Saberes Agroecológicos Tradicionais em Comunidades Quilombolas}


\begin{flushright}
Catuxe Varjão de Santana Oliveira  \footnote[1]{Programa de Pós-graduação em Ciências da Propriedade Intelectual-PPGPI/Universidade Federal de Sergipe-UFS. *E-mail:  \href{mailto:catuxe@academico.ufs.br}{\color{blue}{\underline{catuxe@academico.ufs.br}}}}
\\[0.3cm]
Luiz Diego Vidal Santos \footnote[2]{Universidade Estadual de Feira de Santana - UEFS, Brasil. E-mail: \href{mailto:ldvsantos@uefs.br}{\color{blue}{\underline{ldvsantos@uefs.br}}}}
\\[0.3cm]
Paulo Roberto Gagliardi \footnotemark[1]
\end{flushright}

% resumo em português — bookmark suprimido para manter hierarquia do capítulo
{\renewcommand{\PRIVATEbookmarkthis}[1]{}%
\begin{resumoumacoluna}
A implementação efetiva de mecanismos de valoração bioeconômica para Saberes e Sistemas Agrícolas Tradicionais (SSAT) demanda uma arquitetura de governança que articule os instrumentos técnicos desenvolvidos nos capítulos precedentes --- instrumento WOCAT-SLM-QBR adaptado (Capítulo~2), variáveis linguísticas elicitadas via Delphi (Capítulo~2) e Índice de Valoração Bioeconômica fuzzy (Capítulo~3) --- com os marcos regulatórios de propriedade intelectual coletiva, os mecanismos de repartição de benefícios e as estruturas de tomada de decisão comunitária. Este estudo emprega metodologia de \textit{design participativo} (\textit{co-design}) junto a comunidades quilombolas do Território de Identidade Semiárido Nordeste II (Bahia) para conceber, prototipdar e validar um sistema integrado de governança bioeconômica. O design articula cinco componentes: (i)~plataforma digital de documentação e monitoramento de SSAT baseada no WOCAT-SLM-QBR; (ii)~painel de indicadores (dashboard IVB) com visualização das dimensões de valoração; (iii)~protocolo de certificação comunitária com selo de origem; (iv)~framework jurídico de proteção de conhecimentos tradicionais conforme Lei nº~13.123/2015 e Protocolo de Nagoia; e (v)~modelo de negócio para inserção em cadeias de valor da bioeconomia. A validação seguirá ciclos iterativos de prototipagem, teste em campo e refinamento participativo, assegurando usabilidade, pertinência cultural e sustentabilidade institucional. O estudo contribui para a operacionalização de frameworks de governança que traduzam a riqueza biocultural em mecanismos concretos de empoderamento econômico e proteção jurídica das comunidades detentoras de saberes tradicionais.

\vspace{\onelineskip}
\noindent
\textbf{Palavras-chave}: Design participativo; Governança bioeconômica; Propriedade intelectual coletiva; Sistemas de apoio à decisão; Comunidades quilombolas; Certificação comunitária.
\end{resumoumacoluna}}

\newpage

\section{Introdução}

%% A ser desenvolvido --- OE5: Projetar e validar, mediante design participativo com as comunidades quilombolas,
%% um sistema integrado de governança bioeconômica que articule os instrumentos técnicos
%% (WOCAT-SLM-QBR, Delphi, IVB fuzzy, sinalização de mercado) com marcos regulatórios
%% de PI coletiva, mecanismos de repartição de benefícios e estruturas de tomada de decisão comunitária.

%% Q5: Em que medida um sistema de governança bioeconômica, co-desenhado com as comunidades,
%% articula instrumentos de valoração multidimensional com mecanismos de proteção jurídica
%% e inserção mercadológica de forma culturalmente legítima e institucionalmente sustentável?

%% H5: O design participativo produz arquitetura de governança com maior aderência comunitária,
%% usabilidade percebida e sustentabilidade institucional do que modelos prescritivos
%% concebidos externamente.

\textit{Capítulo em desenvolvimento. Conteúdo a ser elaborado após conclusão dos Capítulos 2, 3 e 4.}


\section{Referencial Teórico}

\textit{A ser desenvolvido.}


\section{Materiais e Métodos}

\textit{A ser desenvolvido.}


\section{Resultados Esperados}

\textit{A ser desenvolvido.}


\section{Discussão Preliminar}

\textit{A ser desenvolvido.}
